% Version: 2026.1.0.0
% ============================================================================
% Engineering Studio AI — Team & Public White Paper
% Author: Hadrian Hu (hadrian.hu@gmail.com)
% Governed by: coding_stds/documentation/latex_style_guide.txt
% Date: \today
% ============================================================================
\documentclass[12pt]{article}

% ============================================================================
% PREAMBLE
% ============================================================================

% --- Font & Encoding ---------------------------------------------------
\usepackage[utf8]{inputenc}
\usepackage[T1]{fontenc}
\usepackage{times}

% --- Page Layout ---------------------------------------------------------
\usepackage[margin=1in]{geometry}
\usepackage{setspace}
\singlespacing

% --- Headers & Footers ----------------------------------------------------
\usepackage{fancyhdr}
\usepackage{lastpage}
\pagestyle{fancy}
\fancyhf{}
\rhead{\thepage\ of \pageref{LastPage}}
\lhead{Engineering Studio AI --- Team \& Public White Paper}
\renewcommand{\headrulewidth}{0.4pt}

% --- Math ------------------------------------------------------------------
\usepackage{amsmath}
\usepackage{amssymb}

% --- Graphics & Tables -------------------------------------------------------
\usepackage{graphicx}
\usepackage{booktabs}
\usepackage{array}
\usepackage{multirow}
\usepackage{float}
\usepackage{tabularx}
\usepackage{ragged2e}

% --- Lists -------------------------------------------------------------------
\usepackage{enumitem}
\setlist[enumerate]{itemsep=2pt, parsep=2pt}
\setlist[enumerate]{itemsep=2pt, parsep=2pt}

% --- Typography safety -------------------------------------------------------
\usepackage[final]{microtype}
\setlength{\emergencystretch}{15pt}
\widowpenalty=10000
\clubpenalty=10000

% --- TOC ------------------------------------------------------------------
\usepackage{tocloft}
\setcounter{tocdepth}{2}

% --- Hyperlinks -----------------------------------------------------------
\usepackage[hidelinks]{hyperref}
\hypersetup{
    colorlinks=false,
    pdftitle={Engineering Studio AI: Team and Public White Paper},
    pdfauthor={Hadrian Hu},
    pdfsubject={Multi-Agent Engineering Workflow Overview}
}

% --- Custom commands ---------------------------------------------------------
\newcommand{\code}[1]{\texttt{#1}}
\newcommand{\FN}[1]{\textsc{fn-#1}}

% --- Long-filename citation macros ------------------------------------------
\newcommand{\VisionClean}{\code{VISION\_MULTIMODAL\_\allowbreak DEPLOYABLE\_\allowbreak AGENTS\_CLEAN.md}}
\newcommand{\VisionGov}{\code{VISION\_AGENTS\_\allowbreak COORDINATION\_\allowbreak ORCHESTRATION\_\allowbreak GOVERNANCE.md}}
\newcommand{\VisionFunc}{\code{VISION\_AGENTS\_\allowbreak BY\_FUNCTION.md}}
\newcommand{\VisionFlow}{\code{VISION\_AGENT\_\allowbreak WORKFLOW\_\allowbreak INTERAGENT\_\allowbreak COORDINATION.md}}
\newcommand{\VisionHackathon}{\code{VISION\_AMD\_\allowbreak LABLAB\_\allowbreak HACKATHON\_\allowbreak ENGINEERING\_STUDIO.md}}

% ============================================================================
\begin{document}

% ============================================================================
% TITLE PAGE
% ============================================================================
\begin{titlepage}
    \centering
    \vspace*{1cm}
    {\LARGE \textbf{Engineering Studio AI}}\\[0.3cm]
    {\large A Multi-Agent, Multi-Modal Platform for the Full}\\
    {\large Software and Hardware Engineering Lifecycle}\\[0.2cm]
    {\large \textit{Team \& Public White Paper}}

    \vspace{2cm}

    {\large \textbf{Author}}\\
    Hadrian Hu\\
    \texttt{hadrian.hu@gmail.com}

    \vspace{1cm}

    {\large CodingStandardsRef Research Notes}\\
    Multimodal Deployable Agents Platform (MDAP)

    \vspace{2cm}

    {\large \today}

    \vfill
    \noindent\rule{\textwidth}{0.4pt}\\[0.2cm]
    {\small
    This document is the team-and-public companion to the formal
    mathematical paper
    \textit{Engineering Studio AI: Formal Model and Mathematical Proofs}
    (v2026.1.0.0).
    The formal paper establishes the same architectural claims as
    rigorous theorems; this document explains the same ideas
    in plain language for a broader audience.
    All architectural claims made here are backed by proofs in the
    companion paper.
    }
\end{titlepage}

% ============================================================================
% ABSTRACT
% ============================================================================
\newpage
\section*{Abstract}

\noindent
Engineering Studio AI is a new kind of software development platform.
Instead of one large AI model trying to do everything at once, it
assembles a coordinated \emph{team} of narrowly focused AI agents ---
specialists, reviewers, testers, adversarial challengers, and a quality
gate --- that collaborate on the full engineering lifecycle from a
single natural-language product brief all the way to a reviewed,
tested, and certified deliverable.

\medskip\noindent
This white paper explains what the platform does, why it is built the
way it is, how its key safety and reliability properties have been
formally verified, and what it means for teams and organizations that
want to adopt AI-assisted engineering at scale.
No mathematical background is required to read this document.
Readers who want the rigorous proofs are referred to the companion
formal paper.

\bigskip
\noindent\textbf{Keywords:}
multi-agent systems,
AI engineering,
workflow automation,
software lifecycle,
hardware emulation,
trust and safety,
MDAP,
Engineering Studio AI

\newpage

% ============================================================================
% TABLE OF CONTENTS
% ============================================================================
\tableofcontents
\newpage

% ============================================================================
\section{What Is Engineering Studio AI?}
\label{sec:what}

\subsection{The One-Sentence Version}

Engineering Studio AI takes a plain-English product brief and
automatically assembles a team of AI specialists that design, review,
challenge, validate, test, document, and certify the resulting
engineering deliverable --- covering both software and emulated
hardware --- without any single agent ever being responsible for
checking its own work.

\subsection{The Problem It Solves}

Building complex engineering products today requires many different
disciplines working in concert: software architecture, firmware,
electrical design, mechanical design, simulation, cost estimation,
legal compliance, testing, and documentation.
Coordinating all of these disciplines is expensive, error-prone, and
slow --- especially when the team is large, distributed, or
working across many simultaneous product lines.

Current AI tools either:
\begin{enumerate}
    \item act as a single, general-purpose assistant that tries to do
          everything but lacks deep discipline expertise, or
    \item require manual orchestration by engineers who must themselves
          know when to call which tool.
\end{enumerate}

Engineering Studio AI takes a different approach: it treats the
engineering process as a \emph{workflow} that can be decomposed,
formalized, and executed by a coordinated team of specialist agents,
each with a narrow, well-defined role, operating under a strict set of
governance rules that prevent any agent from silently certifying its
own output.

\subsection{The Platform in Plain English}

When a user submits a product brief, here is what happens:

\begin{enumerate}
    \item \textbf{An Orchestrator reads the brief.}
          The root Orchestrator agent (\FN{01}) decomposes the brief
          into a set of parallel and sequential tasks, one for each
          relevant engineering discipline.

    \item \textbf{Specialists go to work.}
          Domain Specialist agents (\FN{02}/\FN{03}) --- for example,
          a Firmware Specialist, an Electrical Specialist, a Mechanical
          Specialist --- each work on their assigned piece of the
          deliverable independently.

    \item \textbf{Reviewers check each output.}
          A Reviewer agent (\FN{04}) examines each specialist's
          deliverable against the relevant standard, checklist, or
          specification.  The Reviewer never wrote the deliverable;
          it only critiques it.

    \item \textbf{Adversarial Challengers try to break it.}
          A dedicated Challenge Division agent (\FN{09}) actively
          looks for failure modes, edge cases, and hidden assumptions
          that the Specialist and Reviewer may have missed.
          Its entire job is to ask: ``How does this fail?''

    \item \textbf{A Validator looks for contradictions.}
          The Validator agent (\FN{05}) checks whether the deliverable
          is internally consistent and whether it contradicts any other
          artifact in the system.  Unlike the Reviewer, the Validator
          never fixes anything --- it only reports contradictions.

    \item \textbf{Testers verify behavior.}
          Testing agents (\FN{07}) run the relevant automated tests
          or simulations to confirm that the deliverable behaves as
          specified.

    \item \textbf{Documentation is generated.}
          A Documentation agent (\FN{06}) produces the required
          documentation artefacts alongside the deliverable.

    \item \textbf{A Quality Gate makes the final call.}
          Only the Quality Gate agent (\FN{08}) can issue a formal
          approval certificate.
          If the deliverable passes, it is certified and handed off.
          If it fails, the run loops back for rework --- but only up
          to a fixed, pre-defined number of times, after which the
          system escalates rather than looping forever.
\end{enumerate}

Every one of these roles is performed by a \emph{different} agent.
No agent reviews its own work.
No agent can certify its own output.
This is not a convention --- it is enforced by the architecture.

% ============================================================================
\section{Why This Architecture?}
\label{sec:why}

\subsection{The Problem with a Single AI Agent}

Imagine asking one person to simultaneously design a product, review
their own design, challenge their own assumptions, validate their own
conclusions, run their own tests, and then certify their own work as
production-ready.
No serious engineering organization works this way --- for good
reason.
Self-review is not review.
Self-certification is not certification.

The same principle applies to AI agents.
A single large-context AI model that performs all of these functions on
the same artifact suffers from three independently fatal problems:

\begin{enumerate}
    \item \textbf{Responsibility collapse.}
          When the same agent both builds and reviews a deliverable,
          there is no independent party to catch self-inconsistency.
          The agent is both judge and defendant.

    \item \textbf{Context overload.}
          A monolithic agent must simultaneously hold the full state
          of every engineering discipline, every in-flight artifact,
          all review history, and all governance rules.
          As the project grows, the signal-to-noise ratio of every
          individual decision degrades.

    \item \textbf{Unverifiable trust.}
          If a single agent produces and certifies its own output,
          there is no principled way to assign that output a lower
          trust level than the user's own instructions.
          This creates a subtle but critical safety vulnerability:
          the agent's own conclusions could silently override the
          governance constraints it is supposed to operate within.
\end{enumerate}

Any one of these three problems is sufficient to disqualify a
single-agent architecture for multi-disciplinary production use.
Engineering Studio AI addresses all three by construction.

\subsection{The Adopted Solution: Separation of Concerns at Scale}

Engineering Studio AI applies the same principle that makes large
engineering organizations work:
\emph{separate the people who build from the people who review,
and separate the people who review from the people who certify.}

In the platform, this separation is formalized as the
\textbf{Non-Overlap Rule}: no single agent instance may perform both
a construction role and a critique role on the same artifact.
This rule is not a guideline --- it is enforced structurally, and the
companion formal paper proves that it requires \emph{at least four
distinct agent instances} to cover the mandatory phases of any
artifact lifecycle.

\subsection{Statelessness: Each Agent Knows Only Its Job}

Every agent in the platform is \emph{stateless}: it receives a
well-defined task specification, performs its function, returns its
output, and retains no memory of the interaction.
This design choice has three important consequences:

\begin{enumerate}
    \item \textbf{Predictability.}
          A stateless agent produces the same output given the same
          input.  Its behavior can be tested, audited, and reproduced.

    \item \textbf{Scalability.}
          Stateless agents can be run in parallel without coordination
          overhead.  Six domain specialists can work simultaneously on
          six parallel tasks without needing shared state.

    \item \textbf{Security.}
          A stateless agent cannot accumulate context across sessions,
          so it cannot be gradually manipulated over time into
          violating its governance constraints.
\end{enumerate}

% ============================================================================
\section{The Agent Catalog}
\label{sec:catalog}

\subsection{Twelve Function Classes}

The platform organizes all agent roles into twelve function classes,
labeled \FN{00} through \FN{11}.
Table~\ref{tab:functions} summarizes these classes, their layer in the
organizational hierarchy, and their role in plain English.

\begin{table}[H]
\centering
\caption{The Twelve Agent Function Classes of the MDAP Platform.}
\label{tab:functions}
\begin{tabular}{@{}lllp{7cm}@{}}
\toprule
\textbf{Class} & \textbf{Layer} & \textbf{Name} & \textbf{What it does} \\
\midrule
\FN{00} & Layer 0 & Constitution &
  The non-negotiable governance rules.
  Never produces deliverables; everything else must operate within
  its constraints. \\
\FN{01} & Layer 1 & Orchestrator &
  Reads the brief, decomposes tasks, dispatches agents, and
  merges results. The conductor of the ensemble. \\
\FN{02} & Layer 2 & Domain Specialist &
  Builds the primary deliverable for one engineering discipline
  (e.g., firmware, electrical, mechanical). \\
\FN{03} & Layer 3 & Micro Specialist &
  Handles a narrow sub-task within a domain (e.g., one specific
  component of a PCB layout). \\
\FN{04} & Layer 4 & Reviewer &
  Critiques a deliverable against exactly one standard or checklist.
  Never builds; never fixes. \\
\FN{05} & Layer 5 & Validator &
  Detects contradictions between artifacts or between an artifact
  and its specification.  Never fixes --- only reports. \\
\FN{06} & Layer 6 & Documentation &
  Produces and maintains documentation artifacts alongside the
  engineering deliverable. \\
\FN{07} & Layer 7 & Testing &
  Designs and executes automated tests or simulation runs to
  verify behavior. \\
\FN{08} & Layer 8 & Quality Gate &
  The \emph{sole} agent authorized to issue a formal approval
  certificate.  No other agent can certify a deliverable. \\
\FN{09} & Challenge Div. & Adversarial Challenger &
  Actively attempts to find failure modes, edge cases, and
  hidden assumptions.  Its job is: ``How does this fail?'' \\
\FN{10} & Guardrails & Guardrails &
  Enforces cross-cutting mandates (safety, compliance, ethical
  constraints) at every phase. \\
\FN{11} & Utility & Utility &
  Provides shared, read-only support capabilities (e.g., code
  search, standard lookups) to any agent that needs them. \\
\bottomrule
\end{tabular}
\end{table}

\subsection{The Non-Overlap Rule in Plain English}

The most important governance rule in the platform is simple:
\textbf{the agent that builds a thing cannot be the same agent that
reviews, validates, or adversarially challenges that thing.}

This rule is not a best practice --- it is a hard architectural
constraint.
The companion formal paper proves that, as a direct consequence of
this rule, any valid deployment requires \emph{at least four distinct
agent instances} to handle one artifact's lifecycle:
one to build it, one to review it, one to validate it, and one to
adversarially challenge it.
No deployment with three or fewer agent instances can satisfy this
rule and cover all mandatory phases.

\subsection{The Trust Tier Hierarchy}

Not all instructions and outputs in the system carry the same weight.
The platform assigns every piece of content a \emph{trust tier}:

\begin{enumerate}
    \item \textbf{Tier 0 (T-0):} Model constraints --- the
          non-negotiable safety rules baked into the platform itself.
          Nothing can override these.
    \item \textbf{Tier 1 (T-1):} The Layer~0 constitution --- the
          organization's own governance rules.
    \item \textbf{Tier 2 (T-2):} Verified user instructions.
    \item \textbf{Tier 3 (T-3):} Everything produced by a subagent,
          including all specialist outputs, review reports, and
          validation results.
\end{enumerate}

The key invariant --- proved formally in the companion paper --- is
this: \textbf{a subagent's output starts at Tier~3 and can only be
promoted to a higher trust tier by an explicit Quality Gate
certificate.}
No amount of individual subagent invocations, re-framings, or
chaining can promote a Tier-3 output to Tier-2 status without that
explicit certification step.
This prevents a subtle class of attacks where a sequence of
individually-valid instructions gradually escalates an untrusted
output into a trusted one.

% ============================================================================
\section{The Workflow}
\label{sec:workflow}

\subsection{Eleven Phases, One Deliverable}

Every artifact produced by the platform passes through eleven
workflow phases in a defined order.
The phases are not arbitrary --- they map directly to the natural
structure of an engineering review process.
Table~\ref{tab:phases} lists them in order.

\begin{table}[H]
\centering
\caption{The Eleven Workflow Phases.}
\label{tab:phases}
\begin{tabular}{@{}rll@{}}
\toprule
\textbf{\#} & \textbf{Phase} & \textbf{What happens} \\
\midrule
1  & Plan          & Orchestrator reads brief, defines scope and tasks \\
2  & EnvVerify     & Guardrails agent confirms the environment is safe to proceed \\
3  & StdMap        & Standards applicable to this deliverable are identified \\
4  & Impl          & Domain and Micro Specialists build the deliverable \\
5  & Review        & Reviewer critiques the deliverable (runs concurrently with 6) \\
6  & Challenge     & Adversarial Challenger tries to break the deliverable \\
7  & Validate      & Validator checks for internal contradictions \\
8  & Test          & Testing agent runs automated tests or simulations \\
9  & Doc           & Documentation agent produces required documentation \\
10 & CommitSync    & Changes are committed and synchronized \\
11 & Gate          & Quality Gate issues the final approval or rejection \\
\bottomrule
\end{tabular}
\end{table}

\subsection{Forward Progress and Rework}

The workflow is designed so that phases always move \emph{forward}.
There are no loops in the forward workflow: once the Gate phase
completes, the run is done.

However, the Gate phase can \emph{reject} a deliverable, sending it
back to the beginning of the cycle for rework.
This is the platform's rework mechanism.
Each rework cycle consumes one unit of a pre-configured rework budget.
When the budget is exhausted, the system escalates rather than looping
forever.

The companion formal paper proves that this design guarantees
termination: \textbf{every run of the platform reaches either an
approved state or an escalated state within a finite, pre-computable
number of steps.}
The maximum number of steps is $11 \times (1 + R)$, where $R$ is the
per-phase rework budget configured at deployment time.

\subsection{Concurrent Work and Join Points}

Not all phases are strictly sequential.
The Review phase and the Challenge phase run \emph{concurrently}: the
Reviewer and the Adversarial Challenger both examine the
implementation output in parallel, then the workflow joins at the
Validate phase.

This parallel structure is safe because:
\begin{enumerate}
    \item The two agents are examining the \emph{same} artifact
          independently, not modifying it.
    \item Neither agent can issue a final decision; only the Quality
          Gate can do that.
    \item The workflow only advances past the join point once
          \emph{both} the Review and Challenge phases are complete.
\end{enumerate}

% ============================================================================
\section{Hardware and Multi-Modal Capabilities}
\label{sec:hardware}

\subsection{What ``Multi-Modal'' Means Here}

Engineering Studio AI is described as a multi-modal platform because
it handles inputs and outputs beyond plain text:

\begin{enumerate}
    \item \textbf{Design files:} schematics, PCB layouts, mechanical
          CAD models.
    \item \textbf{Simulation outputs:} SPICE circuit simulation
          results, Gazebo robot simulation outputs.
    \item \textbf{Bill-of-materials:} structured cost and component
          data.
    \item \textbf{Code artifacts:} firmware images, software modules,
          test suites.
    \item \textbf{Documentation:} structured technical documentation,
          compliance reports.
\end{enumerate}

Agents in the platform are not limited to reading and writing text;
they can consume and produce any of these artifact types, depending on
their domain specialization.

\subsection{The Hardware Boundary}

The current instantiation of the platform covers
\emph{simulated and emulated hardware} --- not physical fabrication.
``Hardware'' in this context means:
\begin{enumerate}
    \item Robot simulation in Gazebo.
    \item Circuit simulation in SPICE.
    \item Bill-of-materials generation.
    \item Firmware targeting an emulated processor.
\end{enumerate}

Physical fabrication --- actually manufacturing a PCB, assembling
a robot, or flashing firmware to real silicon --- is outside the
current scope.
This boundary is a disclosed design choice, not a gap:
the platform is designed to validate engineering decisions in
simulation \emph{before} any physical resource is committed.

% ============================================================================
\section{The Formal Guarantees (in Plain English)}
\label{sec:guarantees}

The companion formal paper provides rigorous mathematical proofs of
four core properties of the platform.
This section explains what those proofs mean in practical terms,
without the mathematics.

\subsection{Guarantee 1: The Workflow Always Makes Progress}

The eleven workflow phases form a directed acyclic graph (DAG) ---
a fancy way of saying that the phases are arranged in a strict
forward order with no circular dependencies.

\medskip
\noindent\textbf{What this means in practice:}
There is no combination of phase assignments that can cause the
workflow to spin in a circle without making progress.
Every legitimate execution path through the workflow eventually
reaches the Gate phase.

\medskip
\noindent\textbf{Why it matters:}
Without this guarantee, a subtle ordering bug in the workflow
configuration could cause the system to loop between phases
indefinitely, consuming resources without ever producing a
deliverable.
The formal proof rules this out by construction.

\subsection{Guarantee 2: Rework Always Terminates}

The rework loop --- where a rejected deliverable is sent back for
revision --- is bounded.
Specifically, no run can exceed $11 \times (1 + R)$ total phase
transitions, where $R$ is the deployment-configured per-phase retry
limit.

\medskip
\noindent\textbf{What this means in practice:}
Before a run starts, you can compute the worst-case number of
AI invocations.
If $R = 2$ (two rework cycles per phase), the worst case is
$11 \times 3 = 33$ phase transitions.
No run can exceed this budget; it will either succeed or escalate.

\medskip
\noindent\textbf{Why it matters:}
Without this guarantee, a pathological sequence of rejections could
cause a run to loop indefinitely, making cost estimation and SLA
guarantees impossible.
The formal proof establishes that this scenario cannot occur.

\subsection{Guarantee 3: Responsibilities Never Silently Collapse}

The Non-Overlap Rule requires that the agent who builds a deliverable
is never the same agent who reviews, validates, or adversarially
challenges it.
The formal paper proves this requires \emph{at least four distinct
agents} per artifact lifecycle.

\medskip
\noindent\textbf{What this means in practice:}
A deployment cannot satisfy the governance rules with fewer than four
agents handling the mandatory phases.
Any system configuration that tries to merge the Specialist, Reviewer,
Validator, and Adversarial Challenger roles into fewer than four
distinct agents is provably non-compliant.

\medskip
\noindent\textbf{Why it matters:}
Without this lower bound, a resource-constrained deployment might
silently collapse multiple roles into a single agent, re-introducing
the self-review problem that the architecture was designed to
eliminate.
The formal proof makes this collapse detectable and auditable.

\subsection{Guarantee 4: Trust Levels Cannot Be Forged}

A subagent's output is always treated as Tier-3 (the lowest trust
level) until the Quality Gate explicitly certifies it.
No sequence of subagent invocations --- no matter how many, no matter
how they are chained --- can promote a Tier-3 output to a higher tier
without that explicit certification.

\medskip
\noindent\textbf{What this means in practice:}
An adversarial prompt that tries to convince a subagent to claim
it has higher authority than it actually has cannot succeed.
The trust tier is tracked by the platform, not by the agents
themselves, and can only be changed by the Quality Gate.

\medskip
\noindent\textbf{Why it matters:}
Trust escalation attacks are a real concern in multi-agent systems:
a sequence of individually-valid steps can sometimes lead to a state
that violates a global invariant.
The formal proof shows that this cannot happen in the trust-tier
hierarchy of Engineering Studio AI.

% ============================================================================
\section{The AMD LabLabAI Hackathon Instantiation}
\label{sec:hackathon}

\subsection{Overview}

The first concrete instantiation of Engineering Studio AI was built
for the AMD LabLabAI Hackathon.
This instantiation demonstrates the full platform workflow on a
multi-disciplinary hardware-software engineering problem, covering six
parallel specialist domains: Mechanical, Electrical, Firmware,
Simulation, Cost, and Legal.

\subsection{The Agent Team for One Product Brief}

\begin{table}[H]
\centering
\caption{Agent Roles in the AMD LabLabAI Hackathon Instantiation.}
\label{tab:hackathon-agents}
\begin{tabular}{@{}lll@{}}
\toprule
\textbf{Agent} & \textbf{Function Class} & \textbf{Role} \\
\midrule
Orchestrator           & \FN{01} & Decomposes brief; dispatches specialists \\
Mechanical Specialist  & \FN{02} & Designs mechanical components \\
Electrical Specialist  & \FN{02} & Designs electrical/PCB components \\
Firmware Specialist    & \FN{02} & Writes firmware for the target platform \\
Simulation Specialist  & \FN{02} & Builds and runs simulation models \\
Cost Specialist        & \FN{02} & Produces bill-of-materials and cost model \\
Legal Specialist       & \FN{02} & Checks regulatory and compliance requirements \\
Challenge Division     & \FN{09} & Adversarially reviews all specialist outputs \\
Quality Gate           & \FN{08} & Issues the final certification \\
\bottomrule
\end{tabular}
\end{table}

\subsection{How the Numbers Work Out}

With six domain specialists and one Orchestrator turn per phase,
the worst-case number of subagent invocations for a run with
rework budget $R$ is:
\[
N \;\leq\; 6 \times 11 \times (1 + R) \;=\; 66\,(1 + R).
\]
For a zero-rework run ($R = 0$), this is at most 66 subagent calls.
For one rework cycle ($R = 1$), at most 132 calls.

These are hard upper bounds, not averages.
A typical successful run will use far fewer calls, because not every
phase requires all six specialists and most phases will not require
rework.

\subsection{Hardware Scope}

All hardware work in this instantiation is in simulation or emulation:
\begin{enumerate}
    \item Mechanical: Gazebo robot simulation.
    \item Electrical: SPICE circuit simulation.
    \item Firmware: emulated processor target.
    \item BOM: generated cost model, not physical procurement.
\end{enumerate}
Physical fabrication is explicitly out of scope for this
instantiation.

% ============================================================================
\section{Why This Matters for Teams}
\label{sec:teams}

\subsection{For Engineering Teams}

\begin{enumerate}
    \item \textbf{Every deliverable gets independent review.}
          The architecture ensures that no specialist ever reviews
          their own work.
          For engineering teams, this means every AI-generated artifact
          has been examined by at least three independent agents before
          it reaches the Quality Gate.

    \item \textbf{The system tells you when it cannot succeed.}
          When the rework budget is exhausted, the system escalates
          rather than looping.
          The escalation report tells you exactly which phase failed
          and how many rework cycles were attempted.
          You never get a silent infinite loop.

    \item \textbf{Trust levels are transparent.}
          Every piece of content in the system has a visible trust
          tier.
          A Tier-3 output (specialist-generated) is clearly
          distinguished from a Tier-2 user instruction and from a
          Tier-0 system constraint.
          You always know what has been formally certified and what
          has not.

    \item \textbf{The workflow is extensible without starting over.}
          Adding a new engineering discipline to the platform means
          adding a new Domain Specialist agent with a well-defined
          function class.
          The four formal guarantees (forward progress, bounded rework,
          role separation, trust invariant) automatically extend to
          cover the new agent, provided the Non-Overlap Rule is
          respected.
          No re-proof is required.
\end{enumerate}

\subsection{For Organizations Evaluating AI Adoption}

\begin{enumerate}
    \item \textbf{Auditability by design.}
          Every decision in the workflow is traceable to a specific
          agent with a specific role.
          The audit trail is not an afterthought --- it is a
          structural property of the architecture.

    \item \textbf{Quantifiable worst-case costs.}
          Before a run starts, the maximum number of AI invocations
          is computable from the rework budget and the number of
          phases.
          Cost estimation for AI-assisted engineering is not
          guesswork.

    \item \textbf{No silent capability creep.}
          The Non-Overlap Rule prevents any single agent from
          accumulating responsibilities over time.
          An agent deployed as a Reviewer cannot later be quietly
          reassigned to also produce the artifacts it reviews without
          a deliberate, auditable change to its function-class
          assignment.

    \item \textbf{Governance-first architecture.}
          The platform's governance rules --- trust tiers, rework
          bounds, role separation --- are not operational policies
          bolted on after the fact.
          They are structural properties of the system, verified by
          formal mathematical proof.
          This means they cannot be accidentally disabled by a
          configuration change or a software update that does not
          touch the core architecture.
\end{enumerate}

% ============================================================================
\section{Relationship to the Formal Paper}
\label{sec:formal}

\subsection{What the Formal Paper Proves}

The companion formal paper,
\textit{Engineering Studio AI: Formal Model and Mathematical Proofs}
(v2026.1.0.0), establishes the following four results as rigorous
mathematical theorems:

\begin{enumerate}
    \item \textbf{Theorem~1 (Acyclicity):} The nominal workflow graph
          is a directed acyclic graph.
          This is proved by constructing an explicit labeling function
          that assigns a strict numerical index to each phase, such
          that every phase dependency points from a lower index to a
          higher one --- ruling out any circular dependency by a
          simple arithmetic argument.

    \item \textbf{Theorem~2 (Bounded Termination):} Every run of the
          artifact state machine reaches either an approved or an
          escalated state within at most $P + B$ phase transitions,
          where $P = 11$ is the number of phases and $B = R \times P$
          is the total rework budget.
          This is proved using a potential function (a counter) that
          strictly decreases at every step and is bounded below by
          zero, guaranteeing that the sequence of steps must be finite.

    \item \textbf{Theorem~3 (Minimum Agent Multiplicity):} Any
          Non-Overlap-compliant deployment requires at least four
          distinct agent instances to cover the mandatory phases of
          one artifact lifecycle.
          This is proved by showing that the four mandatory function
          classes (Specialist, Reviewer, Validator, Adversarial
          Challenger) form a complete conflict graph: every pair of
          them conflicts.
          By a standard combinatorial argument (pigeonhole principle),
          this means no compliant deployment can cover all four with
          fewer than four agents.

    \item \textbf{Theorem~4 (Trust-Tier Invariant):} For every
          subagent-originated content item, across an invocation
          sequence of any length, its trust tier remains at Tier~3
          unless the Quality Gate has explicitly certified it.
          This is proved by structural induction: the base case
          is an empty session (trivially true); the inductive step
          examines all three possible event types (user request,
          subagent return, Quality Gate certificate) and shows that
          none of them can promote a Tier-3 item without an explicit
          Quality Gate event.
\end{enumerate}

A fifth result --- a bound on the total number of AI invocations ---
is provided as an explicitly labeled proof sketch.
The sketch is structurally complete, but rests on a per-invocation
cost model that has not yet been empirically validated against a real
deployment.
The formal paper states precisely what additional work is needed to
promote the sketch to a full proof.

\subsection{What This White Paper Does Not Prove}

This white paper explains the architectural ideas and their practical
implications in plain language.
It does not contain mathematical proofs.
All formal claims in this document are backed by the companion formal
paper; readers who require the rigorous justification are referred there.

In particular, the following claims are \emph{not} proved in this
document:
\begin{enumerate}
    \item The semantic correctness of any specialist agent's output
          (e.g., whether a generated PCB layout is electrically valid).
    \item Whether any specific LLM model will produce correct
          specialist-domain outputs.
    \item Cost estimates for any specific deployment configuration.
\end{enumerate}
These are runtime concerns that depend on the specific models,
tooling, and domain specifications used in a deployment.
They are outside the scope of both this document and the formal paper.

% ============================================================================
\section{Frequently Asked Questions}
\label{sec:faq}

\begin{enumerate}

    \item \textbf{Does every run use all eleven phases?}

          Not necessarily.
          The eleven phases represent the full lifecycle of a
          deliverable from planning to certification.
          A simple deliverable may not require all phases to be
          active (for example, a software-only deliverable may not
          need a hardware simulation phase).
          However, the phases that \emph{are} active always execute
          in the defined order, and the Quality Gate is always the
          final step.

    \item \textbf{What happens if an agent produces incorrect output?}

          The Review, Challenge, and Validate phases exist precisely to
          catch incorrect specialist output before it reaches the
          Quality Gate.
          If all three catch it, the deliverable is rejected and sent
          for rework.
          If none of them catch it and the Quality Gate passes it, the
          deliverable is certified as of the standards and specifications
          that the agents were given --- but its \emph{domain-specific}
          correctness (e.g., whether a circuit will actually work in
          the real world) depends on how well the agents' specifications
          matched the real problem.
          This is why human review of the specification inputs is still
          important.

    \item \textbf{Can agents communicate with each other directly?}

          No.
          All inter-agent communication is routed through the
          Orchestrator.
          Agents cannot directly call or message each other.
          This design prevents unauthorized information flows and
          ensures the Orchestrator always has a complete view of the
          workflow state.

    \item \textbf{What is the rework budget and how do I set it?}

          The rework budget $R$ is the maximum number of times any
          single phase can be retried before the system escalates.
          It is a deployment-level configuration parameter.
          Setting $R = 0$ means the first rejection immediately
          escalates (useful for fast-fail scenarios).
          Setting $R = 3$ means each phase can be retried up to three
          times before escalation.
          The total worst-case transitions are $11 \times (1 + R)$;
          use this formula to budget your AI invocations.

    \item \textbf{Can I add new agent roles to the platform?}

          Yes.
          New agent roles are added by assigning them one or more of
          the twelve function classes (\FN{00}--\FN{11}) and updating
          the conflict graph to reflect any new Non-Overlap
          constraints.
          As long as the Non-Overlap Rule is respected, the four
          formal guarantees automatically extend to the new role
          without requiring new proofs.

    \item \textbf{Is the platform limited to software engineering?}

          No.
          The platform covers the full software-and-hardware
          engineering lifecycle, including firmware, mechanical
          design, electrical design, simulation, cost modeling, and
          legal compliance.
          The current instantiation uses emulation and simulation for
          hardware; physical fabrication is a planned future extension.

    \item \textbf{Where can I read the rigorous mathematical proofs?}

          The companion formal paper,
          \textit{Engineering Studio AI: Formal Model and Mathematical
          Proofs} (v2026.1.0.0), contains the complete mathematical
          treatment: precise definitions, full theorem statements,
          step-by-step proofs with all hypotheses explicitly stated,
          and treatment of all boundary conditions and degenerate cases.

\end{enumerate}

% ============================================================================
\section{Roadmap and Future Work}
\label{sec:roadmap}

\begin{enumerate}

    \item \textbf{Empirical validation.}
          Running one complete end-to-end vertical slice --- one full
          product brief through all eleven phases on a real,
          non-trivial engineering deliverable --- would provide
          observational evidence to complement the formal guarantees.
          This is the highest-priority near-term milestone.

    \item \textbf{Completing the token-cost proof.}
          The invocation-count bound (at most $66 \times (1 + R)$
          calls for the six-specialist instantiation) is formally
          proved.
          Converting this to a token-cost bound requires validating
          per-agent token-budget constants against a real deployment.
          Once validated, the proof sketch in the companion paper
          can be promoted to a full proof.

    \item \textbf{Tighter multiplicity bounds.}
          The current formal proof establishes that \emph{at least
          four} distinct agents are required.
          Future work will analyze the full conflict graph to derive
          tighter bounds for specific deployment configurations,
          which will improve resource planning for large deployments.

    \item \textbf{Semantic correctness layer.}
          The current formal model governs the \emph{structure} of the
          workflow (who reviews whom, when, and under what trust level)
          but not the \emph{content} of any deliverable.
          A future extension will introduce domain-specific correctness
          predicates that connect governance-layer guarantees to
          domain-layer correctness, enabling end-to-end formal
          assurances.

    \item \textbf{Physical hardware extension.}
          Extending the platform's hardware coverage from emulation
          to physical fabrication is a planned future milestone,
          requiring a separately validated instantiation of the
          formal model.

    \item \textbf{Expanding the MDAP catalog.}
          The Multimodal Deployable Agents Platform (MDAP) catalog
          currently contains approximately 250 agent roles.
          As the catalog grows, the formal vocabulary of the companion
          paper provides a stable foundation for auditing new roles
          without re-deriving the four core theorems.

\end{enumerate}

% ============================================================================
\section{Conclusion}
\label{sec:conclusion}

Engineering Studio AI represents a principled answer to the question
of how to apply AI effectively to complex, multi-disciplinary
engineering work.

Its core insight is simple: \textbf{the same structural rules that
make large engineering organizations work --- separation of
construction from review, review from certification, and human
oversight of escalation --- apply just as much to AI agent teams as
they do to human ones.}

By formalizing these rules as a workflow graph, an artifact state
machine, a trust-tier hierarchy, and a conflict graph, and by proving
that the resulting system is acyclic, terminating, responsibility-safe,
and trust-invariant, Engineering Studio AI provides something rare in
the current AI landscape: a multi-agent system whose key safety and
reliability properties are not just claimed or empirically observed,
but \emph{mathematically proved}.

The platform is designed to be extended.
New disciplines, new agent roles, new hardware targets can all be
added without invalidating the core guarantees, as long as the
governance rules are respected.
The formal vocabulary of the companion paper is the stable foundation
on which that extension can be built.

\bigskip
\noindent
For the complete mathematical treatment, see:
\medskip

\noindent
Hu, H. (2026).
\textit{Engineering Studio AI: Formal Model and Mathematical Proofs}
(v2026.1.0.0).
CodingStandardsRef Research Notes, Multimodal Deployable Agents Platform.

% ============================================================================
\newpage
\section*{References}
\label{sec:references}

\begin{sloppypar}
\begin{enumerate}
    \item Hu, H. (2026). ``VISION\_MULTIMODAL\_\allowbreak DEPLOYABLE\_AGENTS\_CLEAN.md.''
          Internal repository document, \texttt{markdowns/visions/}.
    \item Hu, H. (2026). ``VISION\_AGENTS\_COORDINATION\_\allowbreak ORCHESTRATION\_\allowbreak GOVERNANCE.md.'' Internal repository document,
          \texttt{markdowns/visions/}.
    \item Hu, H. (2026). ``VISION\_AGENTS\_BY\_FUNCTION.md.'' Internal
          repository document, \texttt{markdowns/visions/}.
    \item Hu, H. (2026). ``VISION\_AGENT\_WORKFLOW\_INTERAGENT\_\allowbreak COORDINATION.md.'' Internal repository document,
          \texttt{markdowns/visions/}.
    \item Hu, H. (2026). ``VISION\_AMD\_LABLAB\_HACKATHON\_ENGINEERING\_\allowbreak STUDIO.md.'' Internal repository document,
          \texttt{markdowns/visions/}.
    \item Hu, H. (2026). \textit{Engineering Studio AI: Formal Model and
          Mathematical Proofs} (v2026.1.0.0). CodingStandardsRef Research
          Notes, Multimodal Deployable Agents Platform.
    \item Cormen, T. H., Leiserson, C. E., Rivest, R. L., \& Stein, C.
          (2009). \textit{Introduction to Algorithms} (3rd ed.). MIT Press.
    \item Rosen, K. H. (2019). \textit{Discrete Mathematics and Its
          Applications} (8th ed.). McGraw-Hill.
\end{enumerate}
\end{sloppypar}

% ============================================================================
\newpage
\section*{Changelog}
\label{sec:changelog}

\begin{table}[H]
\centering
\caption{Document Revision History (YPSV: YEAR.MAJOR.MINOR.PATCH).}
\begin{tabular}{@{}lllp{7.0cm}@{}}
\toprule
\textbf{Version} & \textbf{Date} & \textbf{Author} & \textbf{Description} \\
\midrule
2026.1.0.0 & 2026-07-05 & Hadrian Hu &
  Initial creation: team-and-public white paper companion to the
  formal mathematical paper.
  Covers the architecture, the agent catalog, the workflow, the
  formal guarantees (in plain English), the AMD LabLabAI Hackathon
  instantiation, team implications, FAQ, and roadmap. \\
\bottomrule
\end{tabular}
\end{table}

\end{document}
% ============================================================================
% END OF DOCUMENT
% ============================================================================
